\documentclass[aspectratio=169]{beamer}

\usepackage[utf8]{inputenc}
\mode<presentation>
\usetheme[visual,slim,framecount]{tptng}

\usepackage[T1]{fontenc}
\graphicspath{{images/}{pdfs/}{figs/}{svgs/}{odgs/}}
\DeclareGraphicsExtensions{.pdf,.png,.jpg}

\usepackage[french]{babel}
\usepackage[babel]{microtype}
\usepackage{booktabs}
\usepackage{tikz}
\usetikzlibrary{positioning}
\usepackage{listings}
\usepackage{xcolor}

% XML listing style
\lstdefinelanguage{XML}{
  morestring=[b]",
  morestring=[s]{>}{<},
  morecomment=[s]{<?}{?>},
  stringstyle=\color{teal},
  identifierstyle=\color{blue!70!black},
  keywordstyle=\color{red!70!black},
  morekeywords={root,BehaviorTree,Sequence,Fallback,Action,Condition,
                LoadMission,MissionStructureValid,GenerateInspectionPath,
                NavigateToInspectionPoint,PerformVisualInspection,
                AnalyseMeasurements,Inverter,RetryUntilSuccessful}
}
\lstset{
  basicstyle=\ttfamily\scriptsize,
  breaklines=true,
  frame=none,
  language=XML,
  showstringspaces=false
}

\AtBeginSection[]{
  \contentsframe[currentsection]
}
\renewcommand*{\contentsframename}{Plan}

%------------------------------------------------------------------------------
\title{BTGenBot $\times$ NAV4RAIL}
\subtitle{Génération de Behavior Trees avec LLMs compacts}

\author{Benjamin Lepourtois \and Anne Faury \and Mohamed Amar \and Mourad Latoundji}
\institute{Télécom Paris}
\date[19 mai 2026]{19 mai 2026 \\ {\small Mastère spécialisé IA/DATA — Promotion 2026}}
%------------------------------------------------------------------------------

\begin{document}

\maketitle

%------------------------------------------------------------------------------
\section{Contexte \& Problème}
%------------------------------------------------------------------------------

\begin{frame}[fragile]
  \frametitle{Objectif SNCF : générer des templates de Behavior Trees}

  \begin{columns}
    \begin{column}{0.55\textwidth}
      \textbf{Ce que veut la SNCF :}
      \begin{itemize}
        \item Décrire une mission en langage naturel
        \item Obtenir un \textbf{template BT paramétré} en XML
        \item Template instancié par un fichier de config au runtime
        \item Nœuds = catalogue fermé de \textbf{28 skills} NAV4RAIL
      \end{itemize}

      \vspace{0.5em}
      \textbf{Contrainte principale :}
      \begin{alertblock}{}
        \centering Un seul exemple réel disponible à ce jour
      \end{alertblock}
    \end{column}
    \begin{column}{0.42\textwidth}
      \begin{block}{Template SNCF (cible)}
        \begin{lstlisting}
<Sequence>
  <LoadMission
    mission_file_path=
      "{mission_file_path}"/>
  <MissionStructureValid/>
  <GenerateInspectionPath
    origin_sph="{origin_sph}"
    dest_sph="{dest_sph}"/>
  <Fallback>
    <NavigateToInspectionPoint/>
    <HandleFailure/>
  </Fallback>
</Sequence>
        \end{lstlisting}
      \end{block}
    \end{column}
  \end{columns}
\end{frame}

%------------------------------------------------------------------------------

\begin{frame}
  \frametitle{Pourquoi on ne peut pas générer le dataset proxy ?}

  \begin{columns}
    \begin{column}{0.48\textwidth}
      \begin{alertblock}{Tentative dataset proxy}
        \begin{itemize}
          \item Génération synthétique par LLM
          \item Résultat : \textbf{très loin des BTs réels}
          \item Structure des BTs incorrecte, câblage des ports erroné
          \item \textbf{Aucune connaissance métier ferroviaire}
        \end{itemize}
      \end{alertblock}

      \vspace{0.4em}
      \begin{exampleblock}{Solution retenue}
        \small Reproduire un article de référence sur \textbf{données publiques réelles},
        puis argumenter le transfert SNCF.
      \end{exampleblock}
    \end{column}
    \begin{column}{0.48\textwidth}
      \textbf{Ce qu'il faudrait idéalement :}
      \begin{itemize}
        \item 20–50 templates BT réels SNCF
        \item Descriptions associées en NL
        \item Fichier de config avec variables blackboard
      \end{itemize}

      \vspace{0.5em}
      \textbf{Ce qu'on a :}
      \begin{itemize}
        \item 1 exemple réel (mission d'inspection)
        \item Catalogue de 28 skills avec ports
        \item 27 règles de sécurité (L1–L5)
      \end{itemize}
    \end{column}
  \end{columns}
\end{frame}

%------------------------------------------------------------------------------
\section{BTGenBot — Article de référence}
%------------------------------------------------------------------------------

\begin{frame}
  \frametitle{BTGenBot : Politecnico di Milano (2024)}
  \framesubtitle{R.A. Izzo — github.com/AIRLab-POLIMI/BTGenBot}

  \begin{columns}
    \begin{column}{0.5\textwidth}
      \textbf{Hypothèse centrale :}
      \begin{block}{}
        \small Un LLM compact (7B paramètres) fine-tuné sur \textbf{600 exemples réels}
        génère des BTs de qualité comparable aux grands modèles.
      \end{block}

      \vspace{0.3em}
      \textbf{4 contributions :}
      \begin{enumerate}\small
        \item Dataset 600 paires BT~$\leftrightarrow$~description NL
        \item Système d'évaluation sur 9 tâches
        \item Comparaison 3 LLMs $\times$ base vs fine-tuné
        \item Pipeline exécution sur robot physique
      \end{enumerate}
    \end{column}
    \begin{column}{0.46\textwidth}
      \textbf{Pipeline BTGenBot :}

      \vspace{0.3em}
      \begin{tikzpicture}[
        box/.style={draw, rounded corners, fill=blue!10, minimum width=2.8cm,
                    minimum height=0.6cm, font=\small, text centered},
        arr/.style={->, thick, blue!60}
      ]
        \node[box] (nl) at (0,3) {Description NL};
        \node[box] (llm) at (0,2) {LLM fine-tuné (7B)};
        \node[box] (xml) at (0,1) {BT en XML};
        \node[box] (ros) at (0,0) {Exécution ROS2};
        \draw[arr] (nl) -- (llm);
        \draw[arr] (llm) -- (xml);
        \draw[arr] (xml) -- (ros);
        \node[draw, dashed, red!60, right=0.2cm of ros, font=\tiny,
              minimum width=1.4cm, text width=1.4cm, text centered]
              {hors scope\\reproduction};
      \end{tikzpicture}
    \end{column}
  \end{columns}
\end{frame}

%------------------------------------------------------------------------------

\begin{frame}
  \frametitle{Méthode BTGenBot — Dataset \& Fine-tuning}

  \begin{columns}
    \begin{column}{0.48\textwidth}
      \textbf{Construction du dataset :}
      \begin{enumerate}
        \item \textbf{600 BTs réels} extraits de projets ROS2 open-source
              (Nav2, manipulation, exploration) — validés sur robots
        \item \textbf{GPT-3.5-turbo} génère la description NL de chaque BT
              (prompt one-shot, max 200 mots)
        \item Format \textbf{Alpaca} : \texttt{instruction} / \texttt{input} / \texttt{output}
      \end{enumerate}
    \end{column}
    \begin{column}{0.48\textwidth}
      \textbf{Fine-tuning LoRA :}

      {\scriptsize
      \begin{tabular}{ll}
        \toprule
        Paramètre & Valeur \\
        \midrule
        Modèles & LLaMA-2-7b / Chat / CodeLlama \\
        Méthode & LoRA (r=8, $\alpha$=16) \\
        Modules & q/k/v/o\_proj + MLP (7) \\
        Epochs & 70–80 \\
        LR & 3e-4 \\
        Contexte & 2048 tokens \\
        GPU & 2$\times$ RTX Quadro 6000 \\
        \bottomrule
      \end{tabular}}
    \end{column}
  \end{columns}
\end{frame}

%------------------------------------------------------------------------------

\begin{frame}
  \frametitle{Résultats BTGenBot — Chiffres clés}

  \begin{columns}[T]
    \begin{column}{0.36\textwidth}
      \textbf{\small Baseline zero-shot}\\
      {\tiny (grands modèles, sans FT)}
      \vspace{0.3em}

      {\scriptsize
      \begin{tabular}{p{2.4cm}cc}
        \toprule
        Modèle & Synt. & Sém. \\
        \midrule
        GPT-3.5 (175B+) & 100\% & 77.8\% \\
        Gemini ($\sim$1T) & 88.9\% & 55.6\% \\
        LLaMA-2-13b & 33.3\% & 22.2\% \\
        \midrule
        \textbf{LLaMA-2-7b} & \textbf{0\%} & — \\
        \bottomrule
      \end{tabular}}

      \vspace{0.3em}
      {\tiny Les 7B sans FT sont inutilisables.}
    \end{column}

    \begin{column}{0.30\textwidth}
      \textbf{\small Après fine-tuning (ZS)}\\
      {\tiny Phase 1 — tâches 1–7}
      \vspace{0.3em}

      {\scriptsize
      \begin{tabular}{lcc}
        \toprule
        Modèle & Synt. & Sém. \\
        \midrule
        LLaMAChat-7b FT & 71.4\% & 57.1\% \\
        CodeLlama-7b FT & \textbf{85.7\%} & \textbf{100\%} \\
        \bottomrule
      \end{tabular}}
    \end{column}

    \begin{column}{0.30\textwidth}
      \textbf{\small Validation complète}\\
      {\tiny Phase 2 — one-shot + correction}
      \vspace{0.3em}

      {\scriptsize
      \begin{tabular}{lc}
        \toprule
        Config. & BT OK \\
        \midrule
        LLaMAChat-7b OS+SA & \textbf{77.7\%} \\
        CodeLlama-7b OS+SA & 33.3\% \\
        \bottomrule
      \end{tabular}}

      \vspace{0.3em}
      {\tiny OS = one-shot \quad SA = static analysis}
    \end{column}
  \end{columns}

  \vspace{0.5em}
  \begin{alertblock}{}
    \centering\small
    \textbf{Nouvelle baseline à évaluer :} Claude Opus 4.7 (Anthropic, mai 2026)
    en zero-shot sur les 9 tâches — quelle est la limite des grands modèles sans fine-tuning en 2026 ?
  \end{alertblock}
\end{frame}

%------------------------------------------------------------------------------
\section{Lien avec la SNCF}
%------------------------------------------------------------------------------

\begin{frame}
  \frametitle{De BTGenBot aux templates SNCF}

  \begin{columns}
    \begin{column}{0.48\textwidth}
      \begin{block}{Ce que BTGenBot prouve}
        {\small\begin{itemize}
          \item 600 exemples + LoRA $\Rightarrow$ BTs valides
          \item Domaine ouvert (ROS2 varié)
          \item Nœuds libres, ports non typés
          \item BTs statiques (pas de templates)
        \end{itemize}}
      \end{block}
    \end{column}
    \begin{column}{0.48\textwidth}
      \begin{exampleblock}{SNCF : domaine plus contraint}
        {\small\begin{itemize}
          \item \textbf{Vocabulaire fermé} : 28 skills fixes
          \item \textbf{Ports typés} par le catalogue
          \item \textbf{Variables blackboard} : \texttt{\{mission\_fp\}}
          \item \textbf{Templates} instanciés par config
        \end{itemize}}
      \end{exampleblock}
    \end{column}
  \end{columns}

  \vspace{0.3em}
  \begin{alertblock}{\centering Argument central}
    \centering\small
    Domaine contraint $=$ espace de génération réduit $=$ \textbf{moins d'exemples nécessaires}.\\[0.1em]
    \textbf{20–50 templates réels SNCF suffisent} pour reproduire les performances de BTGenBot.
  \end{alertblock}
\end{frame}

%------------------------------------------------------------------------------

\begin{frame}
  \frametitle{Catalogue-conditioned generation : la clé du transfert SNCF}

  \begin{columns}
    \begin{column}{0.50\textwidth}
      \textbf{Différence architecturale avec BTGenBot :}

      \vspace{0.5em}
      Le catalogue est injecté \textbf{dans le prompt} à chaque inférence :

      \vspace{0.3em}
      \begin{tikzpicture}[
        box/.style={draw, rounded corners, minimum height=0.55cm,
                    font=\small, text centered},
        arr/.style={->, thick}
      ]
        \node[box, fill=orange!20, minimum width=2.8cm] (cat) at (0,2.5) {Catalogue 28 skills};
        \node[box, fill=blue!10,   minimum width=2.8cm] (nl)  at (0,1.5) {Mission NL};
        \node[box, fill=green!10,  minimum width=2.8cm] (llm) at (0,0.5) {LLM fine-tuné};
        \node[box, fill=teal!10,   minimum width=2.8cm] (tpl) at (0,-0.4){Template XML};
        \draw[arr] (cat) -- (llm);
        \draw[arr] (nl)  -- (llm);
        \draw[arr] (llm) -- (tpl);
      \end{tikzpicture}
    \end{column}
    \begin{column}{0.46\textwidth}
      \textbf{\small Ce que le LLM apprend avec 20–50 exemples :}
      {\small\begin{itemize}
        \item Sélectionner les skills pertinents parmi 28
        \item Câbler les ports avec les bonnes \texttt{\{variables\}}
        \item Respecter les prérequis (\texttt{LoadMission} avant \texttt{Valid})
        \item Produire un template réutilisable
      \end{itemize}}

      \vspace{0.2em}
      \begin{block}{}
        \small Le catalogue agit comme \textbf{prior structurel fort} :
        le LLM choisit parmi 28 options, il n'invente plus.
      \end{block}
    \end{column}
  \end{columns}
\end{frame}

%------------------------------------------------------------------------------
\section{Plan de reproduction}
%------------------------------------------------------------------------------

\begin{frame}
  \frametitle{Ce qu'on reproduit — et ce qu'on adapte}

  \begin{columns}[T]
    \begin{column}{0.48\textwidth}
      \textbf{Reproduit directement :}

      \vspace{0.2em}
      {\scriptsize\begin{tabular}{p{2.8cm}p{3cm}}
        \toprule
        Composant & Source \\
        \midrule
        Dataset & 600 paires (GitHub public) \\
        Fine-tuning LoRA & QLoRA, modèles 2024 \\
        Éval. syntaxique & Validateur XML/BT \\
        Éval. sémantique & TED + Node-F1 + LLM-judge \\
        Base vs FT & Tableau 9 catégories \\
        \bottomrule
      \end{tabular}}
    \end{column}
    \begin{column}{0.48\textwidth}
      \textbf{Adapté / non reproduit :}

      \vspace{0.2em}
      {\scriptsize\begin{tabular}{p{2.4cm}p{2.4cm}}
        \toprule
        Composant & Raison \\
        \midrule
        Simulation ROS2 & Non disponible \\
        Robot physique & Hors scope \\
        LLaMA-2 original & $\to$ LLaMA-3.1/Qwen2.5 \\
        \bottomrule
      \end{tabular}}

      \vspace{0.4em}
      \begin{exampleblock}{Évaluation sémantique sans ROS2}
        {\scriptsize\begin{itemize}
          \item \textbf{TED} : nb d'opérations pour transformer le BT généré en BT référence
          \item \textbf{Node-F1} : quels nœuds du BT référence sont présents dans la génération
          \item \textbf{LLM-judge} : Phi-3-mini dit si le BT correspond à la description NL (77.8\% accord GPT-4, validé dans l'article)
        \end{itemize}}
      \end{exampleblock}
    \end{column}
  \end{columns}
\end{frame}

%------------------------------------------------------------------------------

\begin{frame}
  \frametitle{Plan de travail — 4 semaines}

  {\scriptsize
  \begin{columns}[T]
    \begin{column}{0.50\textwidth}
      \begin{block}{Semaine 1 — Appropriation du dataset}
        \begin{itemize}\itemsep0pt
          \item Téléchargement dataset BTGenBot (GitHub)
          \item \textbf{EDA} : catégories, longueur BTs, vocabulaire nœuds
          \item Conversion JSONL + adaptation prompt
        \end{itemize}
      \end{block}

      \begin{block}{Semaine 2 — Fine-tuning}
        \begin{itemize}\itemsep0pt
          \item SFT QLoRA sur 3 modèles (cluster Télécom, RTX 3090)
          \item LLaMA-3.1-8B, Qwen2.5-7B, Qwen2.5-Coder-7B
          \item $\sim$2–4h par modèle
        \end{itemize}
      \end{block}
    \end{column}
    \begin{column}{0.46\textwidth}
      \begin{block}{Semaine 3 — Évaluation}
        \begin{itemize}\itemsep0pt
          \item Correction syntaxique (validateur BT)
          \item TED + Node-F1 vs référence
          \item LLM-judge sur 9 catégories
          \item Zero-shot vs one-shot
        \end{itemize}
      \end{block}

      \begin{block}{Semaine 4 — Analyse \& Rédaction}
        \begin{itemize}\itemsep0pt
          \item Tableau comparatif base vs FT
          \item Discussion transfert SNCF
          \item Rédaction rapport final
        \end{itemize}
      \end{block}
    \end{column}
  \end{columns}}
\end{frame}

%------------------------------------------------------------------------------
\section{Conclusion}
%------------------------------------------------------------------------------

\begin{frame}
  \frametitle{Synthèse — BTGenBot comme premier entonnoir}

  \begin{columns}[T]
    \begin{column}{0.50\textwidth}

      \vspace{0.2em}
      \begin{tikzpicture}[
        box/.style={draw, rounded corners, font=\small, text centered,
                    minimum height=0.7cm, minimum width=4.0cm},
        arr/.style={->, thick, gray}
      ]
        \node[box, fill=gray!15] (gen) at (0, 3.2) {LLM généraliste (7B)};
        \node[box, fill=blue!15] (bt)  at (0, 2.1)
          {\textbf{BTGenBot} {\scriptsize — 600 ex. ROS2}};
        \node[box, fill=orange!20] (sncf) at (0, 0.8)
          {\textbf{Modèle SNCF} {\scriptsize — fine-tuning SNCF}};
        \node[box, fill=green!20] (out) at (0, -0.3)
          {Templates BT NAV4RAIL};
        \draw[arr] (gen)  -- (bt);
        \draw[arr] (bt)   -- (sncf);
        \draw[arr] (sncf) -- (out);
        % entonnoir 1 icon (between gen and bt)
        \node[font=\normalsize, text=blue!60]  at (0.55, 2.72) {$\triangledown$};
        \node[font=\tiny,       text=blue!60]  at (0.55, 2.50) {1er entonnoir};
        % entonnoir 2 icon (between bt and sncf)
        \node[font=\normalsize, text=orange!70] at (0.55, 1.48) {$\triangledown$};
        \node[font=\tiny,       text=orange!70] at (0.55, 1.26) {2e entonnoir};
      \end{tikzpicture}
    \end{column}
    \begin{column}{0.46\textwidth}
      \textbf{\small Ce que BTGenBot apporte :}
      {\small\begin{itemize}\itemsep2pt
        \item Prouve qu'un LLM 7B peut apprendre la \textbf{syntaxe et la sémantique des BTs}
        \item Fournit un modèle \textbf{déjà spécialisé} — pas un LLM généraliste en point de départ
        \item Réduit l'effort du second fine-tuning SNCF
      \end{itemize}}

      \vspace{0.1em}
      \begin{alertblock}{Message clé}
        {\scriptsize BTGenBot = \textbf{1\textsuperscript{er} entonnoir} (domaine BT).
        Le 2\textsuperscript{e} entonnoir (domaine SNCF) reste à construire
        dès que la SNCF fournit ses exemples réels.}
      \end{alertblock}
    \end{column}
  \end{columns}
\end{frame}

%------------------------------------------------------------------------------
\appendix
%------------------------------------------------------------------------------

\begin{frame}
  \frametitle{Annexe — Hyperparamètres de reproduction}

  \begin{center}
  {\small\begin{tabular}{lp{3.6cm}p{4.0cm}}
    \toprule
    \textbf{Paramètre} & \textbf{BTGenBot original} & \textbf{Notre reproduction} \\
    \midrule
    Modèles & LLaMA-2-7b, Chat, CodeLlama & LLaMA-3.1-8B, Qwen2.5-7B/Coder \\
    PEFT & LoRA & QLoRA (NF4) \\
    \texttt{lora\_r} & 8 & 8–16 \\
    \texttt{lora\_alpha} & 16 & 16–32 \\
    Modules cibles & 7 (q/k/v/o + MLP) & identique \\
    Epochs & 70–80 & 70–80 \\
    Learning rate & 3e-4 & 2e-4–3e-4 \\
    Dataset & 600 paires & 600 paires (identique) \\
    GPU & 2$\times$ RTX Quadro 6000 & RTX 3090 (cluster Télécom) \\
    \bottomrule
  \end{tabular}}
  \end{center}
\end{frame}

%------------------------------------------------------------------------------

\begin{frame}
  \frametitle{Annexe — 9 tâches d'évaluation BTGenBot}

  \begin{columns}
    \begin{column}{0.48\textwidth}
      \begin{enumerate}
        \item \textbf{Navigation} : visiter des positions dans l'ordre
        \item \textbf{Navigation avec priorité} : seuil sur lecture capteur
        \item \textbf{Navigation avec fallback} : sauter waypoints inaccessibles
        \item \textbf{Navigation + bras} : activer manipulateur à chaque position
        \item \textbf{Exploration} : navigation continue, check complétion
      \end{enumerate}
    \end{column}
    \begin{column}{0.48\textwidth}
      \begin{enumerate}
        \setcounter{enumi}{5}
        \item \textbf{Exploration manipulateur} : cycler configs articulaires
        \item \textbf{Vision active + saisie} : observer, estimer prise, pick-and-place
        \item \textbf{Traitement matériaux} : séquence de boutons + évaluation statut
        \item \textbf{Assemblage multi-station} : collecter composants, assembler
      \end{enumerate}
    \end{column}
  \end{columns}

  \vspace{0.5em}
  \small Phase 1 (sélection) = tâches 1–7 \quad Phase 2 (validation) = tâches 1–9
\end{frame}

\end{document}
